\chapter{系统需求分析}

\section{系统建设目标}
本系统面向人工智能专业的医学图像处理实验，目标是搭建一个能实际操作的 MRI 医学图像拟生成平台。平台需要完成数据预处理、StyleGAN2 二维切片生成、连续切片组织、点云初始化、3DGS 训练结果展示和 AI 辅助说明等工作。和只运行命令行脚本相比，平台更强调“能看见过程、能找到结果、能复核路径”。

使用对象主要包括三类：进行医学图像生成实验的学生，学习 StyleGAN2 与 3DGS 的学习者，以及需要展示生成式医学影像流程的教学人员。平台不面向临床诊断，也不输出医学诊断结论，其定位是算法教学、毕业设计展示和工程实验验证。

\section{功能需求分析}
平台功能围绕一条实际使用路线展开：先处理 MRI 数据，再生成二维切片，最后查看三维结果。用户进入平台后，应能看到数据、模型、任务和结果资产，而不是只看到脚本执行后的文件夹。

数据预处理是第一步。用户提供 NIfTI 数据目录后，平台需要调用预处理脚本，完成轴向切片、灰度归一化、尺寸缩放和低信息切片过滤。处理结束后，系统要返回输出路径、有效切片数量和日志信息。如果路径不存在、文件格式不对或没有生成有效切片，页面应给出明确提示。

模型训练功能用于得到二维生成权重。训练任务耗时较长，因此不能让浏览器一直等待。后端需要在后台启动训练进程，并保存配置文件、训练日志和模型快照。前端只负责提交参数和查看状态。这样处理后，用户可以知道训练是否启动、结果保存在哪里，也能在论文复核时找到对应证据。

切片生成功能把已经训练好的 \texttt{.pkl} 权重变成可查看图片。用户设置随机种子、生成数量和截断参数后，后端加载权重并输出 PNG 图像。生成结果既保存在结果目录中，也会返回给前端显示。这个功能是平台展示 StyleGAN2 效果的主要入口。

连续序列和点云初始化用于连接二维与三维。系统通过潜变量插值得到一组连续切片，再把非背景像素转换为空间点，最后写成 PLY 文件。这样得到的点云不是临床三维模型，而是一个工程中间结果，用来说明二维切片怎样进入后续 3DGS 流程。

3DGS 结果管理负责保存和展示三维训练产物。平台需要能查看训练日志、相机参数、checkpoint、渲染图像和训练后点云。三维模块的价值不只是“生成一个文件”，还在于让这些文件可以被追踪、解释和再次打开。

AI 辅助问答主要服务于教学说明。它可以解释 MRI 切片生成、StyleGAN2 参数、3DGS 概念和平台操作步骤，但不参与核心训练和推理。即使外部 AI 接口不可用，数据处理、图像生成、点云初始化和三维结果展示仍应保持可用。

\section{性能与可用性需求分析}
系统性能需求主要来自耗时任务。MRI 批量预处理、StyleGAN2 训练和 3DGS 训练都可能持续较长时间，因此应由后端任务进程执行。React 前端只负责提交任务、展示日志和刷新状态。切片生成、图像预览、影像分析和点云统计更接近演示操作，应尽量在较短时间内返回结果。

可用性方面，平台需要减少用户直接操作命令行的次数。常用参数通过输入框、滑块和按钮填写，运行日志在页面中显示。路径错误、权重文件缺失、生成数量异常或三维资产不存在时，系统应说明问题出在哪里，而不是只返回程序异常。正常运行时，页面应给出图像预览、输出路径和关键日志，方便用户判断结果是否正确。

扩展性主要体现在三个方向。二维生成部分应支持更换 StyleGAN2 权重或训练数据。三维部分应支持继续接入不同切片序列、点云和 3DGS 输出结果。评价部分可以在现有可视化、L1 和 PSNR 基础上继续补充 SSIM、FID、LPIPS、专家评价或下游任务指标。平台代码应保持模块化，避免每增加一个功能就重写整体界面。

\section{可行性分析}
从技术可行性看，本文所用技术都有较成熟的开源实现或 Python、TypeScript 生态支持。StyleGAN2-ADA 有可复用的 PyTorch 工程，NIfTI 数据可由 nibabel 读取，图像处理可使用 NumPy 和 PIL 完成。FastAPI 可以组织后端接口和 WebSocket 通信，React 可以构建复杂前端页面，Vite 可以支撑本地开发和前端构建。当前项目已经形成 StyleGAN2 权重、生成切片、初始点云、3DGS 训练日志和渲染结果，这说明核心流程已经具备实现基础。

从数据可行性看，IXI 数据集公开提供多种脑部 MRI 数据，类型包括 T1、T2、PD、MRA 和 DTI 等。这些数据以 NIfTI 格式发布，适合教学和算法实验\cite{ixiDataset}。本文使用 T1 脑部 MRI 作为主要实验数据，能够支撑二维切片生成和三维序列组织。

从工程可行性看，系统采用“算法脚本 + FastAPI 后端 + React 前端”的组织方式。底层脚本负责计算任务，后端负责接口封装、任务调度和结果归档，前端负责参数输入、状态展示、影像查看和三维交互。三者通过 HTTP 接口、WebSocket、文件路径和日志连接。这样的组织方式更接近真实 Web 平台结构，也便于继续扩展页面和接口。

从应用可行性看，平台定位为教学与研究演示系统，并不是临床诊断软件。平台不承担医疗器械级别的验证要求。论文对生成结果的描述以算法展示、工程闭环和可视化验证为主，不把平台输出解释为临床可用数据。
